\chapter{R\'esoudre des probl\`emes}\label{ch:g2:problems}

Une histoire avec des nombres cache une question~: additionner~?
enlever~? multiplier~? Ce chapitre donne un moyen de d\'ecider,
utilise de simples barres pour \emph{voir} l'op\'eration, et lit des
informations dans des tableaux et des pictogrammes.

\section{Choisir l'op\'eration}

\begin{method}[Quelle op\'eration~?]\label{met:g2:problems:choose}
\begin{enumerate}
  \item \emph{mettre ensemble, en tout}~: additionner ($+$)~;
  \item \emph{enlever, combien restent, combien de plus}~:
        soustraire ($-$)~;
  \item \emph{autant de groupes \'egaux}~: multiplier ($\times$)~;
  \item dessiner un croquis en cas de doute, et toujours r\'epondre
        par une phrase.
\end{enumerate}
\end{method}

\begin{omfigure}
\begin{tikzpicture}[scale=0.6]
  % put together
  \draw[thick, fill=omDef!20] (0,0) rectangle (3,0.8);
  \draw[thick, fill=omThm!20] (3,0) rectangle (4.6,0.8);
  \node[font=\small] at (1.5,0.4) {$24$};
  \node[font=\small] at (3.8,0.4) {$15$};
  \draw[Stealth-Stealth] (0,-0.4) -- (4.6,-0.4);
  \node[below, font=\small] at (2.3,-0.5) {? en tout~: $24 + 15$};
  % compare
  \begin{scope}[xshift=8cm]
  \draw[thick, fill=omDef!20] (0,0.9) rectangle (4.6,1.7);
  \node[font=\small] at (2.3,1.3) {$40$};
  \draw[thick, fill=omThm!20] (0,0) rectangle (2.8,0.8);
  \node[font=\small] at (1.4,0.4) {$28$};
  \draw[Stealth-Stealth] (2.8,0.4) -- (4.6,0.4);
  \node[below, font=\small] at (2.3,-0.2) {? de plus~: $40 - 28$};
  \end{scope}
\end{tikzpicture}

{\small Images en barres~: deux barres coll\'ees signifient une
addition~; deux barres compar\'ees signifient une soustraction.}
\end{omfigure}

\begin{example}[Un probl\`eme en deux \'etapes]\label{ex:g2:problems:twostep}
\guillemotleft{} Sam ach\`ete $3$ paquets de $6$ autocollants, puis
en donne $5$. Combien lui en reste-t-il~? \guillemotright{}
\begin{enumerate}
  \item Groupes~: $3 \times 6 = 18$ autocollants~;
  \item donner~: $18 - 5 = 13$.
\end{enumerate}
Sam garde $13$ autocollants. Relire la question~: l'histoire est-elle
vraiment finie~? Ici oui.
\end{example}

\begin{example}[Un nombre en trop, cach\'e]\label{ex:g2:problems:extra}
\guillemotleft{} Un livre de $210$ pages\,\dots \guillemotright{} ---
attention, tous les nombres d'une histoire ne sont pas n\'ecessaires.
\guillemotleft{} Ben a $8$ billes rouges et $6$ billes bleues et a
$7$~ans. Combien de billes~? \guillemotright{} L'\^age est un
leurre~: $8 + 6 = 14$ billes. Trier les nombres avant de calculer.
\end{example}

\section{Tableaux et pictogrammes}

\begin{example}[Un tableau]\label{ex:g2:problems:table}
G\^ateaux vendus \`a un stand~:
\begin{center}
\renewcommand{\arraystretch}{1.3}
\begin{tabular}{l|ccc}
 & samedi & dimanche & total \\
\hline
muffins & $12$ & $9$ & $21$ \\
tartes & $7$ & $10$ & $17$ \\
\end{tabular}
\end{center}
Lecture~: tartes du dimanche, $10$. Calcul~: muffins sur les deux
jours, $12 + 9 = 21$~; plus de tartes le dimanche que le samedi de
$10 - 7 = 3$.
\end{example}

\begin{example}[Un pictogramme]\label{ex:g2:problems:pictogram}
Lire d'abord la l\'egende~: ici une pomme repr\'esente $2$ fruits.

\begin{omfigure}
\begin{tikzpicture}[scale=0.62]
  \node[left, font=\small] at (0,2) {rouges};
  \foreach \i in {0,1,2,3} \fill[omThm] (0.5+\i*0.75,2) circle (0.24);
  \node[left, font=\small] at (0,1) {vertes};
  \foreach \i in {0,1} \fill[omMeth] (0.5+\i*0.75,1) circle (0.24);
  \node[left, font=\small] at (0,0) {jaunes};
  \foreach \i in {0,1,2} \fill[omProp] (0.5+\i*0.75,0) circle (0.24);
  \node[right, font=\small] at (4.0,1) {l\'egende~: un point $= 2$ fruits};
\end{tikzpicture}

{\small Un pictogramme~: le rouge montre $4$ points, donc
$4 \times 2 = 8$ fruits rouges~; le vert $2$ points, $4$ fruits~; le
jaune $3$ points, $6$ fruits.}
\end{omfigure}
\end{example}

\section{Exercices}

\begin{exercise}[$\star$]\label{exo:g2:problems:1}
Nomme l'op\'eration (sans calculer)~: \guillemotleft{} $34$ et $28$
billes mises ensemble \guillemotright{}~; \guillemotleft{} $50$
bonbons, $12$ mang\'es \guillemotright{}~; \guillemotleft{} $4$
bo\^ites de $5$ pommes \guillemotright{}~; \guillemotleft{} Ana a
$30$, L\'eo a $18$~: combien de plus a Ana~? \guillemotright{}.
\end{exercise}

\begin{exercise}[$\star$]\label{exo:g2:problems:2}
R\'esous~: \guillemotleft{} Un parking a $156$ voitures~; $48$
partent. Combien restent-elles~? \guillemotright{}
\end{exercise}

\begin{exercise}[$\star$]\label{exo:g2:problems:3}
R\'esous~: \guillemotleft{} $5$ enfants cueillent chacun $4$ fleurs.
Combien de fleurs en tout~? \guillemotright{}
\end{exercise}

\begin{exercise}[$\star$]\label{exo:g2:problems:4}
Dessine une image en barres, puis r\'esous~: \guillemotleft{} Tom a
$45$ cartes, Mia en a $17$ de moins. Combien en a Mia~?
\guillemotright{}
\end{exercise}

\begin{exercise}[$\star$]\label{exo:g2:problems:5}
Deux \'etapes~: \guillemotleft{} Une caisse contient $6$ bouteilles.
Un magasin re\c{c}oit $4$ caisses et en vend $9$. Combien de
bouteilles restent-elles~? \guillemotright{}
\end{exercise}

\begin{exercise}[$\star$]\label{exo:g2:problems:6}
Trouve et barre le nombre inutile, puis r\'esous~:
\guillemotleft{} Ben a $8$~ans et $23$ billes. Il en gagne $15$ de
plus. Combien de billes maintenant~? \guillemotright{}
\end{exercise}

\begin{exercise}[$\star$]\label{exo:g2:problems:7}
Avec le tableau de \cref{ex:g2:problems:table}~: combien de tartes
vendues sur les deux jours~? Combien de g\^ateaux de toutes sortes
le samedi~?
\end{exercise}

\begin{exercise}[$\star$]\label{exo:g2:problems:8}
Avec le pictogramme de \cref{ex:g2:problems:pictogram}~: combien de
fruits en tout~? Quelle couleur en a le plus~?
\end{exercise}

\begin{exercise}[$\star$]\label{exo:g2:problems:9}
R\'esous~: \guillemotleft{} Un ruban mesure $80$~cm. L\'ea en coupe
$35$~cm. Quelle longueur reste-t-il~? \guillemotright{}
\end{exercise}

\begin{exercise}[$\star\star$]\label{exo:g2:problems:10}
\guillemotleft{} Un boulanger fait $5$ plateaux de $4$ g\^ateaux le
matin et $3$ plateaux de $4$ g\^ateaux l'apr\`es-midi. Combien de
g\^ateaux ce jour-l\`a~? \guillemotright{} (Deux multiplications,
puis une \omterm{def:g1:addition:def}{somme} --- ou compter d'abord les plateaux~!)
\end{exercise}
